\subsection{Topological properties of the image of a path}

PROPOSITION 15. --- The image of a path in a Hausdorff topological space is a compact, metrizable, connected, and locally path-connected topological space.

Let X be a Hausdorff topological space and \(c\colon I \to X\) a continuous map. Denote by R the equivalence relation \(c(s) = c(t)\) in I. The topological space \(c(\mathbf{I})\) is Hausdorff (TG, I, p. 63, cor. 1), the space I/R is quasi-compact (TG, I, p. 62, th. 2), hence the bijection I/R \(\to c(\mathbf{I})\) deduced from c is a homeomorphism (TG, I, p. 63, cor. 2). Consequently, the space \(c(\mathbf{I})\) is compact, metrizable (TG, IX, p. 22, prop. 17), connected (TG, I, p. 82, prop. 6), and locally path-connected (III, p. 261, prop. 8).

THEOREM 1 (Hahn and Mazurkiewicz). --- Every metrizable, compact, nonempty, connected, and locally connected topological space is homeomorphic to a quotient space of the interval {[}0,1{]}.

Lemma 4. --- Let K be a compact topological space and let R be a set of open sets of K covering K. There exists an entourage V of the uniform structure of K (TG, II, p. 27, th. 1) such that, for every \(x \in K\) , \(V(x)\) is contained in one of the sets belonging to R.

For every point x of K, there exists an entourage \(W_{x}\) of the uniform structure of K such that \(\mathrm{W}_{x}(x)\) is contained in one of the sets belonging to R. Let \(V_{x}\) be an entourage of the uniform structure of K such that \(\overset{2}{\mathrm{V}}_{x}\) is contained in \(W_{x}\). The interiors of the \(\mathrm{V}_{x}(x)\) cover K; since the space K is compact, there exists a finite subset F of K such that the family \((\mathrm{V}_{y}(y))_{y\in\mathrm{F}}\) is a covering of K (TG, I, p. 59). Denote by V the intersection of the family \((\mathrm{V}_{y})_{y\in\mathrm{F}}\); the set V is an entourage of the uniform structure of K. For every point \(x\in K\), there exists a point \(y\in F\) such that x belongs to \(\mathrm{V}_{y}(y)\). Consequently, the set \(\mathrm{V}(x)\) is contained in \(\overset{2}{\mathrm{V}}_{y}(y)\), hence in one of the sets belonging to R.

Lemma 5. --- Let X and Y be uniform spaces, and let f be a mapping from X into Y. Let F be a collection of closed subsets of X covering X and having the following properties:

\begin{enumerate}
\def\labelenumi{(\roman{enumi})}
\item
  There exists a set \(F_{0} \in F\) which meets all the sets \(F \in F\) ;
\item
  For every entourage U of the uniform structure of X, there are only finitely many sets \(F \in F\) which are not small of order U;
\item
  For every entourage V of the uniform structure of Y, there are only finitely many sets \(F \in F\) such that \(f(F)\) is not small of order V.
\end{enumerate}

Then, if the restriction of f to each of the sets \(F \in F\) is continuous, f is continuous.

Let x be a point of X. Let us prove that f is continuous at x. There exists a set \(F_{1} \in F\) such that \(x \in F_{1}\). The restriction of \(f\) to \(F_{0} \cup F_{1}\) is continuous (TG, I, p. 19, prop. 4). Replacing \(F_{0}\) by \(F_{0} \cup F_{1}\), we reduce to the case where \(x \in F_{0}\).

Let V be an entourage of the uniform structure of Y. Choose an entourage \(V'\) of this same uniform structure such that \(V' \subset V\). Since the restriction of f to \(F_0\) is continuous, there exists an entourage U of the uniform structure of X such that \(f(z) \in \mathrm{V}'(f(x))\) for all \(z \in \mathrm{F}_0 \cap \mathrm{U}(x)\). Let \(U'\) be an entourage of the uniform structure of X such that \(U' \subset U\). Let A denote the union of \(F_0\), of the sets \(F \in F\) which are not small of order \(U'\) and of those such that \(f(\mathrm{F})\) is not small of order \(V'\). By hypothesis, A is the union of a finite number of sets belonging to F, and the restriction of f to A is continuous (loc. cit.). There therefore exists a neighborhood W of x in X, contained in \(U'(x)\), such that \(f(y) \in \mathrm{V}(f(x))\) for \(y \in A \cap W\). To conclude, it will suffice to prove that we also have \(f(y) \in \mathrm{V}(f(x))\) for every point \(y \in (\mathrm{X} - \mathrm{A}) \cap \mathrm{W}\). Let y be such a point. Let F be an element of F such that \(y \in F\). By definition of A, F is small of order \(U'\) and \(f(\mathrm{F})\) is small of order \(V'\). By hypothesis, F meets \(F_0\). Let \(z \in F \cap F_0\). We have \(z \in U'(y)\) since F is small of order \(U'\) and \(y \in U'(x)\) since W is contained in \(U'(x)\), whence \(z \in U'(x)\) and a fortiori \(z \in U(x)\). But then, since z belongs to \(F_0\), we have \(f(z) \in V'(f(x))\). Moreover \(f(\mathrm{F})\) is small of order \(V'\), whence \(f(y) \in V'(f(z))\). It follows that one has \(f(y) \in V'(f(x))\) and finally \(f(y) \in V(f(x))\). This concludes the proof of Lemma 5.

Let us now prove Theorem 1. Let X be a nonempty compact, connected, and locally connected metric space. Such a space is path-connected and locally path-connected (III, p. 267, Corollary 2). By the corollary and the example of III, p. 270, there exists a continuous surjective map f from the Cantor ternary set K (TG, IV, p. 9, example) onto X. We shall construct a continuous extension g of f to {[}0,1{]}, which will prove Theorem 1.

Let \(\varepsilon\) be a real number \(>0\). The open, path-connected subsets of X of diameter \(\leqslant \varepsilon\) cover X. Let \(\mathcal{R}\) denote the set of their inverse images under \(f\); this is a set of open subsets of K covering K. By Lemma 4 of III, p. 272, there exists a real number \(\alpha > 0\) such that every closed ball in K of radius \(\alpha\) is contained in an element of \(\mathcal{R}\). In particular, if \(t\) and \(t'\) are points of K such that \(|t - t'| \leqslant \alpha\), there exists a path in X joining \(f(t)\) to \(f(t')\) whose image has diameter \(\leqslant \varepsilon\).

The preceding paragraph allows one to construct by induction a strictly increasing sequence \((n_{k})_{k\geqslant0}\) of integers \(\geqslant0\) having the following property: for every integer \(k\geqslant0\) and every pair \((t,t')\) of points of K such that \(|t-t'| \leqslant3^{-n_{k}}\), there exists a path in X joining \(f(t)\) to \(f(t')\) whose image has diameter \(\leqslant2^{-k}\). The complement of K in {[}0,1{]} is the union of a family \((\mathrm{I}_{n,p})\) of pairwise disjoint open intervals, where n ranges over the set of integers \(\geqslant0\) and p over the set of integers between 1 and \(2^{n}\) (TG, IV, p. 9, example). Consider one of these intervals \(I_{n,p}\) and write it as {]}a,b{[}. The points a and b belong to K and one has \(b-a=3^{-n-1}\). One defines the function g on the interval \(I_{n,p}\) in the following way: one chooses a path c in X joining \(f(a)\) to \(f(b)\), whose image has diameter \(\leqslant2^{-k}\) if \(n_{k}\leqslant n<n_{k+1}\), and one sets \(g(t)=c(\frac{t-a}{b-a})\) for \(t\in I_{n,p}\). The function \(g\colon[0,1]\to X\) thus defined extends f. It is continuous on K as well as on each of the closed intervals \(\overline{I_{n,p}}\). These latter meet K. Moreover, for every real number \(\varepsilon>0\), there are only finitely many intervals \(\overline{I_{n,p}}\) of length \(>\varepsilon\), and there are only finitely many intervals \(\overline{I_{n,p}}\) whose images under g have diameter \(>\varepsilon\). By Lemma 5, the map g is continuous.
% semantic-fragments: legacy-input-seam
\relax{}
